% ============================================================================
\clearpage
\section*{v4.10.0-r48 Source-Only: The Readout of OP1 (Reading B) and of the OP5 Comparison Model Compared}
\addcontentsline{toc}{section}{v4.10.0-r48 Source-Only: Two Readout Notions Compared}
\label{sec:v4100-r48-readout-compare}
% ============================================================================

\paragraph{Purpose.}
OP1 in reading B (the register width is an observer quantity) and the OP5 comparison model on its readout side (v4.10.0-r46, r47) both place
structure in the readout.  This layer checks whether they use the same notion of readout.  It is an exploration with criteria fixed before the
comparison; it is partly not blind, since both definitions were read beforehand.  Earlier layers are not rewritten.

\noindent\textbf{Source-only status.}
This v4.10.0-r48 layer is source-only.  It does not build a PDF, does not assign a new DOI,
and does not authorize a public release.  The v5.3.0 Version DOI
\href{https://doi.org/10.5281/zenodo.24434403}{10.5281/zenodo.24434403} and the Concept DOI
\href{https://doi.org/10.5281/zenodo.19210728}{10.5281/zenodo.19210728} are unchanged by it.

\subsection*{v4.10.0-r48 the two definitions}
\label{subsec:v4100-r48-definitions}

In reading B of OP1 the register measures \(\varphi=|s|/(|\mathbf E|^2+|\mathbf B|^2)\) of the field seen by an observer, with \(s\) the complex invariant; R is the frame with
\(\mathbf E\parallel\mathbf B\) (\(\varphi=1\)), P the null field (\(\varphi=0\)), and M lies at the rapidity \(\eta\) with \(\tanh\eta=\beta_{\rm sym}\).  In the OP5 comparison model the directed
coupling is \(-h\,\Pi^{\mu\nu}u^\alpha u^\beta R_{\mu\alpha\nu\beta}\), the readout of the tidal field transverse to the measuring axis \(n\) of a detector with rest frame \(u\) (r46.1).

\subsection*{v4.10.0-r48 results}
\label{subsec:v4100-r48-results}

\paragraph{Theoremlet v4.10.0-r48.1 (dependence on the observer).}
For a field with \(\mathbf E\parallel\mathbf B\) along \(z\), an observer moving across the axis with rapidity \(\eta\) measures \(|\mathbf E'|^2+|\mathbf B'|^2=(E_0^2+B_0^2)\cosh2\eta\), which is twice
the energy density \(|T_{\mu\nu}u^\mu u^\nu|\), and \(\varphi=\operatorname{sech}2\eta\); an observer moving along the axis sees \(\varphi=1\).  The OP5 numerator of r46.3 is, with \(\beta=\tanh\eta\),
\(\tfrac1{12}[8a_\perp^2+4a_\perp k_z^2+3k_z^4-4k_z^2(a_\perp+k_z^2)\cosh2\eta+k_z^4\cosh4\eta]\) for motion along the axis and a polynomial in \(\cosh2\eta\), \(\cosh4\eta\) for motion
across it.  All of these are even in \(\eta\).
\emph{Proof.} Exact Lorentz transformation of \(F\) and series-free evaluation of r46.3 (script \texttt{readout\_\allowbreak compare.wl}).  \(\square\)

\paragraph{Remark v4.10.0-r48.2 (comparison).}
Against the criteria fixed beforehand the two notions are \emph{related, not the same}.  Common to both: the readout takes the electric part of a
field strength relative to the four-velocity of the observer (\(F\) for OP1, the Riemann tensor for OP5), it is quadratic in \(u\), and so the time
orientation drops out in both; in OP1 only the sign of \(\operatorname{Re}s\) separates M from T.  Different: in OP1 the rest frame R and its axis
(\(\mathbf E\parallel\mathbf B\)) belong to the field and only the rapidity to the observer, and the readout is a dimensionless ratio; in the comparison model the axis
belongs to the detector, the velocity refers to the preparation, and the readout is an absolute tidal quantity.  For the two to coincide, the
measuring axis \(n\) would have to be the \(\mathbf E\parallel\mathbf B\) axis of the field rest frame and the detector velocity would have to refer to R; R would then also be a
candidate for the field-side frame that r47 requires for the slowed propagation.  This assignment is named, not derived or tested.  The two
notions do not contradict each other.

\subsection*{v4.10.0-r48 frames}
\label{subsec:v4100-r48-frames}

R is the rest frame of the field of convention F, M the matter frame; all quantities are compared in the frame of the observer or detector.  Time
reversal and the tachyon map are not used.

\subsection*{v4.10.0-r48 register entries}
\label{subsec:v4100-r48-register}

\begin{center}
\small
\begin{tabular}{@{}>{\raggedright\arraybackslash}p{0.34\textwidth}>{\raggedright\arraybackslash}p{0.58\textwidth}@{}}
\toprule
statement & status after v4.10.0-r48\\
\midrule
readout notions of OP1 (B) and OP5 & related (electric part relative to \(u\), even in \(\eta\)), not the same (axis and rest frame)\\
field rest frame R as frame of the slowing & named candidate, not tested\\
OP1, OP5 & open\\
\bottomrule
\end{tabular}
\end{center}

\subsection*{v4.10.0-r48 notes}
\label{subsec:v4100-r48-notes}

\texttt{notes/\allowbreak research/\allowbreak op1\_\allowbreak op5\_\allowbreak readout\_\allowbreak compare\_\allowbreak 2026\_\allowbreak 11\_\allowbreak 23/} (method note commit 498f0578,
script \texttt{readout\_\allowbreak compare.wl}, output in \texttt{run/}).

\subsection*{v4.10.0-r48 ledger and guard update}
\label{subsec:v4100-r48-ledger-guard}

\[
R_{\rm chk}^{v4100\text{-}r48}=(1,1,0)
\]
for (status recorded; computations checked in the cited script; the assignment of axis and rest frame not derived).

The layer spends no guard:
\[
R_{\rm glob}^{(v4100\text{-}r48,+)}=(0,1,1,1,1),
\qquad
\Gamma_{\rm guard}^{v4100\text{-}r48}=(1,1,1,1,1),
\qquad
\beta_4=1/256\ \text{unspent}.
\]
The layer does not imply exact finite \(C5\) zero closure, exact all-readout terminality, \(C6/T2\), a global photon-clock/RH theorem,
or global all-readout closure.  OP1--OP5, C5 and C6/T2 remain open.
